\paragraph{Author notation and average-case conventions (Section 3).}
For each $n\in\mathbb N$, let $[n]=\{1,\ldots,n\}$. For each $x\in\{0,1\}^n$ and $i\in[n]$, we let $x_i$ denote the $i$-th bit of $x$ and $x_{\leq i}$ denote $x_1\circ\cdots\circ x_i$. For a distribution $D$, we write $x\leftarrow D$ to refer to a random sampling $x$ according to $D$. For a finite set $S$, we also use the notation $x\leftarrow_u S$ to denote the uniform sampling from $S$. For each $n\in\mathbb N$, we let $U_n$ denote the uniform distribution over $\{0,1\}^n$ or a random variable selected uniformly at random from $\{0,1\}^n$ in context. We use the notation $\operatorname{negl}$ to represent a certain negligible function, i.e., for any polynomial $p(n)$, $\operatorname{negl}(n)<1/p(n)$ for sufficiently large $n\in\mathbb N$. For a randomized algorithm $A$ using $r(n)$ random bits on an $n$-bit input, we use $A(x;s)$ to refer to the execution of $A(x)$ with a random tape $s$ for $x\in\{0,1\}^n$ and $s\in\{0,1\}^{r(n)}$.
For any oracles $O_0$ and $O_1$, we let $O_0+O_1$ denote the combination, i.e., for any $b\in\{0,1\}$ and any $x\in\{0,1\}^*$, $(O_0+O_1)(b\circ x)=O_b(x)$.
In this paper, we assume the basic knowledge of probability theory, including the union bound, Markov's inequality, Hoeffding's inequality, and the Borel--Cantelli lemma.
For each $p\in[0,1]$ and set $S$ of variables taking binary values, we define a $p$-random restriction $\rho$ to $S$ as a partial assignment $\rho:S\to\{0,1,*\}$ (where $*$ represents ``unassigned'') randomly selected as follows: for each $x\in S$,
\[\rho(x)=\begin{cases}*&\text{with probability }p,\\0&\text{with probability }(1-p)/2,\\1&\text{with probability }(1-p)/2.\end{cases}\]
For every restriction $\rho$ to $S$ and function $f$ defined on $S$, we let $f|_\rho$ denote the restricted function obtained by applying a partial assignment to $f$ according to $\rho$.
\paragraph{3.1. Average-Case Complexity.}
We present the notions in average-case complexity theory. Further backgrounds can be found in a survey [7].
We say that a family $D=\{D_n\}_{n\in\mathbb N}$ of distributions, where each $D_n$ is a distribution on $\{0,1\}^n$, is (polynomial-time) samplable if there exists a randomized sampling algorithm $S$ such that the distribution of $S(1^n)$ is identical to $D_n$ for each $n\in\mathbb N$. We consider a family of distributions as a single distribution on instances. We define a distributional problem as a pair of a language $L\subseteq\{0,1\}^*$ and a distribution $D=\{D_n\}_{n\in\mathbb N}$ on instances. For a standard complexity class $C$ (e.g., $\mathsf{NP}$), we define its average-case extension $\mathsf{Dist}C$ as $\mathsf{Dist}C=\{(L,D):L\in C,D\text{ is samplable}\}$.
We present the errorless average-case easiness. We say that a distributional problem $(L,D)$ has an errorless heuristic algorithm $A$ with failure probability $\epsilon:\mathbb N\to(0,1)$ if (1) $A$ outputs $L(x)(:=\mathbf1\{x\in L\})$ or $\bot$ (which represents ``failure'') for every $x\in\operatorname{supp}(D)$, and (2) the failure probability that $A(x)$ outputs $\bot$ over the choice of $x\leftarrow D$ is bounded above by $\epsilon(n)$ for each $n\in\mathbb N$. Note that an errorless heuristic algorithm never outputs an incorrect value $\neg L(x)$ for any $x\in\operatorname{supp}(D)$. Then, for every $\epsilon:\mathbb N\to(0,1)$, we define a class $\mathsf{Avg}_\epsilon\mathsf P$ as a class of distributional problems that have a polynomial-time errorless heuristic algorithm with failure probability $\epsilon(n)$. Furthermore, we say that a distributional problem $(L,D)$ has an errorless heuristic scheme $A$ if $A$ is given an instance $x\in\operatorname{supp}(D)$ and $\epsilon\in(0,1)$ as input and satisfies the condition of an errorless heuristic algorithm with failure probability $\epsilon$. We define a class $\mathsf{AvgP}$ as a class of distributional problems that have a polynomial-time errorless heuristic scheme. It is not hard to verify that $\mathsf{AvgP}\subseteq\mathsf{Avg}_{1/p(n)}\mathsf P$ for any polynomial $p(n)$.
Next, we present the error-prone average-case easiness. We say that a distributional problem $(L,D)$ has an error-prone heuristic algorithm $A$ with failure probability $\epsilon:\mathbb N\to(0,1)$ if the failure probability that $A(x)\neq L(x)$ over the choice of $x\leftarrow D$ is bounded above by $\epsilon(n)$ for each $n\in\mathbb N$. Note that an error-prone heuristic algorithm may output an incorrect value $\neg L(x)$, but the error probability is bounded above by $\epsilon(n)$. Then, for every $\epsilon:\mathbb N\to(0,1)$, we define a class $\mathsf{Heur}_\epsilon\mathsf P$ as a class of distributional problems that have a polynomial-time error-prone heuristic algorithm with failure probability $\epsilon(n)$. We also define an error-prone heuristic scheme and the class $\mathsf{HeurP}$ in the same manner as the errorless case.
We also define classes $\mathsf{AvgP/poly}$, $\mathsf{HeurP/poly}$, $\mathsf{AvgSIZE}[s(n)]$, and $\mathsf{HeurSIZE}[s(n)]$ for each size parameter $s(n)$ in the same manner as above.
\paragraph{Definition 11 (auxiliary-input pseudorandom generator).}
Let $n,m:\mathbb N\to\mathbb N$ be polynomials. We say that $G=\{G_z:\{0,1\}^{n(|z|)}\to\{0,1\}^{m(|z|)}\}_{z\in\{0,1\}^*}$ is an auxiliary-input pseudorandom generator (AIPRG) against $C$ if each $G_z(x)$ is polynomial-time computable from $(z,x)$, $n(\ell)<m(\ell)$ holds for any $\ell\in\mathbb N$, and for any adversary $A$ in $C$, there exists an infinite subset $Z_A\subseteq\{0,1\}^*$ such that for every $z\in Z_A$,
\[\left|\Pr[A(z,G_z(U_{n(|z|)}))=1]-\Pr[A(z,U_{m(|z|)})=1]\right|<\operatorname{negl}(|z|).\]
When the auxiliary-input $z$ is obvious in context, we write $n(|z|)$ and $m(|z|)$ as $n$ and $m$, respectively.
\paragraph{Definition 13 (hitting set generator).}
Let $\ell,m:\mathbb N\to\mathbb N$ be polynomials. We say that $G=\{G_n\}_{n\in\mathbb N}$, where $G_n:\{0,1\}^{\ell(n)}\to\{0,1\}^{m(n)}$ is a hitting set generator (HSG) against $C$ if each $G_n$ is polynomial-time computable, $\ell(n)<m(n)$ holds for each $n\in\mathbb N$, and $G$ hits any language recognized by adversaries in $C$ in the following sense: For any adversary $A$ in $C$, let $L_A=\{L_{A,n}\}_{n\in\mathbb N}$ be a language recognized by $A$, where $L_{A,n}\subseteq\{0,1\}^n$ for each $n\in\mathbb N$. Then, for infinitely many $n\in\mathbb N$, the following holds:
\[|L_{A,m(n)}|>2^{m(n)-1}\ \Longrightarrow\ L_{A,m(n)}\cap\operatorname{Im}G_n\neq\varnothing.\]
